\section{Filesystem Path Architecture}

A file operation always begins with a reference to a location. That reference is the path. It may look like a simple string, but it encodes filesystem structure: directory names, final file names, suffixes, roots, separators, and the difference between relative and absolute addressing.

This section explains how Python represents and manipulates paths before any actual file content is read or written. It contrasts plain string paths with \texttt{pathlib.Path} objects, shows how the current working directory affects relative paths, and introduces portable path construction across Windows and POSIX-style systems.

The central rule is simple: a path is not the file itself. It is a structured reference that Python passes to the operating system when the program wants to inspect, create, open, move, or delete filesystem entries.